\section{Represent a larger integer with little-endian limbs}
A \emph{limb} is one fixed-width part of a larger integer. Let \(\beta=2^{32}\), and
let \(a_i\) denote the unsigned word at byte offset \(4i\). The four-word object means
\[
A=a_0+a_1\beta+a_2\beta^2+a_3\beta^3,\qquad 0\leq a_i<\beta.
\]
The index increases with significance and with the address. Hexadecimal numbers on
paper are usually written most significant digit first, so their visual order is the
reverse of this word order. For example, the integer \(2\beta-1\), written as
\texttt{00000001ffffffff}, has limb zero \texttt{ffffffff}, limb one
\texttt{00000001}, and two zero high limbs.

\begin{figure}[htbp]
\centering
\begin{tikzpicture}[alt={A four-limb unsigned integer stored at increasing byte offsets. Limb zero at offset zero has weight one; limbs one, two and three at offsets four, eight and twelve have weights beta, beta squared and beta cubed. Arithmetic carry travels toward increasing offsets.}]
\foreach \i/\offset/\weight in {0/0/1,1/4/{\beta},2/8/{\beta^2},3/12/{\beta^3}} {
  \draw (3*\i,0) rectangle +(2.7,1.2);
  \node at (3*\i+1.35,.8) {\(a_{\i}\)};
  \node at (3*\i+1.35,.3) {weight \(\weight\)};
  \node at (3*\i+1.35,-.4) {offset \offset};
}
\draw[->] (.3,1.6)--(11.7,1.6) node[midway,above]{carry and increasing significance};
\end{tikzpicture}
\caption{Four 32-bit limbs form one unsigned 128-bit integer.}\label{rsa-limb-layout}
\end{figure}

Use a four-byte stride to visit caller objects. Four-byte alignment does not promise
that an object is aligned for \texttt{ld}. Copy its words individually into an
explicitly eight-byte-aligned private stack area before using doubleword loads there.
Keep the distinction between the object and its pointer: adding four to an address
uses \texttt{daddiu}, regardless of the width of the stored limb.

Before multiplying, write a comparison that begins at limb three and moves toward
limb zero. The first unequal pair determines the order. Compare each pair unsigned;
sign-extending a limb and comparing it as a signed word would put \texttt{ffffffff} below zero. If all four pairs are
equal, the objects represent the same integer even when their addresses differ.
